\subsection{Public Transport}
The buildings of MFF are based at several locations in Prague so one of the common things to do for every student is to travel from one building to another. 
The average distance travelled is increased by regularly visiting the Charles University Sports Center (SCUK) in Hostivař because its location is not actually in a comfortable distance for any student of the Charles University.
Anyone can profit from using the Prague public transport because it is one of the best in the world -- and if you use it correctly, you might even save an hour or two a day. You will usually save two hours a day if you ride a bike. 
Although it is only possible to become perfectly time effective through practice, we offer you a few tips on how to reduce the risk of getting lost.

All of the information about timetables, fares and so on can be found on the website of the Prague Public Transit \url{dpp.cz}.

You might also benefit from using smartphone apps -- there is a lot of them so you can choose whichever you find most suitable for you. (Idos, PID Lítačka or Google Maps)


\subsubsection{Fares}
As the students of Matfyz usually travel by the public transport more than once a day, buying individual tickets would not be really smart budget-wise so they usually buy prepaid time tickets. 
Even if you used the 30-day time ticket for only a week it would be more cost-effective than buying the individual tickets. The time ticket can be either electronic or paper.

The 30-day time ticket for students costs 130 CZK, the 90-day time ticket costs 360 CZK and the annual one costs 1280 CZK. 
You can buy them at the Main Control Centre at Na bojišti 5, Prague 2 or at another sales locations \url{www.dpp.cz/prodejni-mista-jizdnich-dokladu/}.

It is possible to use your student discount while buying time tickets all year round (including holidays). The only exception is that if you buy a time ticket at the end of the academic year (or at the beginning of a new one), you will have to show a new student certificate.

\subsubsubsection{Paper time ticket}
If you want to get a paper time ticket, you will need a student ID card (valid ISIC or a special form of the Prague Public Transit filled out by the school).
It is necessary to always have your student ID card or ISIC with you while travelling by the public transport; otherwise the paper time ticket will not be valid.

\subsubsubsection{Electronic time ticket}
The electronic time ticket can be uploaded to Lítačka card (which can also be used as a library card for the Municipal Library of Prague), In Karta card (from the company ČD) or to a contactless credit card. 
Just as in the case of the paper time ticket you should always have your student certificate or ISIC with you -- but actually if you have the electronic time ticket instead of the paper one, most of the times no one really checks this.
\\
\textbf{Buying electronic time tickets}
\\
The first one needs to be bought at one of the Prague Public Transit's sales locations and it is necessary to show your student ID card. 
The other ones can be also bought at the sales locations or on the internet. 
It is also necessary to show a new student certificate at any of the sales locations every August, September or October because it is impossible to buy the time tickets with student discount online until then.


\subsubsubsection{Transit Inspection}
If an inspector stops you, do not try to run away. There are usually more of them and they cooperate with the police. 
If the inspector catches you without any ticket, you will get a penalty of 1500 CZK; if you pay it right away or manage to pay it in the time set by the Terms and Conditions of Carriage (``within 15 days after the date of inspection''), you will pay only 800 CZK. If you just forgot your time ticket at home and manage to bring it in time (between 12:30 pm the following working day and the 15th day from the imposition of the penalty) to the Main Control Centre at Na bojišti, you will pay only 50 CZK. 


\subsubsection{Types of Public Transport}
We usually use metro, tram, bus, funicular, ferry, train, bike or kick scooter to get around Prague. Provided we have a valid time ticket we do not even have to pay any penalties.

\subsubsubsection{Metro}
Metro is the most important part of the whole system. It opens from 5 a.m. until midnight, or more precisely a bit later as midnight is the time the last trains set out from the terminal station.
A lot of time is spent by standing on the escalator. There is also a certain escalator etiquette saying that you are supposed to stand on the right side so other people could walk (or run) on the other side.
You will later find out that travelling by metro can be optimized -- the waiting time in the station of departure might be used to stand in the place where you will need to get off later. 
Some of the students try to make this system perfect by moving inside the train in between the stations.

Metro is the best way to travel long distances. If your trip is only short and a tram is also going in the same direction, it is better to go by tram. 
Travelling by the Prague metro is comfortable even in the highest and lowest temperatures.

\subsubsubsection{Tram}
Trams are perfect for travelling shorter distances. 
Their speed is highly dependent on the current traffic situation as during rush hour the trams are not even able go through in some sections of the route because the cars are blocking the way (especially on the route between stations Malostranská and Újezd).
Every line has its regular route which unfortunately often changes due to traffic closures so you need to frequently check \url{dpp.cz} where you can find everything about the current traffic situation.
Information about the transit restrictions can also be find on information boards in the metro stations or on small yellow boards above the timetables at the bus and tram stops. 
The same yellow boards usually hang in the buses and trams too.

The daytime tram numbers range from 1 to 26 although not all of them are used.

\subsubsubsection{Bus}
If you find yourself in a place where no metro or tram line can be found, you can usually still go by bus. Their routes often run radially from the metro stations (especially the terminal ones).
The buses also can be really crowded. Their numbers range from 100 to 670.
The bus lines from 1xx to 2xx are part of the Prague public transport system and you usually do not need to show your ticket to anyone but, of course, sometimes you can meet an inspector even there.
There are also suburban transport lines (numbers 3xx, 4xx and 6xx) where you get on the bus via the front door and you have to prove to the driver that your ticket is valid. 
If you have an electronic time ticket you can do this by putting it on the card reader and than the driver can easily check out whether you time ticket is valid or not.
Sometimes you do not even need to prove the validity of your ticket -- when suburban buses go in the direction of Prague, you can usually just get on the bus via any door and do not need to show your ticket to anyone unless you meet an inspector.

\subsubsubsection{Trains}
Trains are also a part of the Prague Integrated Transport (Czech abbreviation is PID) which means that you can also use your time ticket to travel by them, although not all of them. 
This applies only to the lines S and R -- most of the regional and regional fast trains running in Prague and some of the fast trains.

One of the routes that is especially interesting for the students of Matfyz is the route from Hlavní nádraží (``the Main Train Station") to Benešov (the route number is 221) because it can be used to travel to Hostivař.
While travelling to Hostivař by metro or bus takes an hour, travelling by train takes only thirteen minutes. The only disadvantage is that the train usually runs just once every thirty minutes (the intervals can vary depending on the time of the day).
It is also necessary to have in mind that the trains to Hostivař depart from the furthest platforms of Hlavní nádraží and it usually takes at least two minutes to run there from the metro station.

The rules for paying penalties of the Prague Public Transit do not apply while travelling by trains. 
Which means that there is no discount for forgetful people -- if you forget your ticket, the best thing you can do is to confess to the conductor and pay the price set by České dráhy + 40 CZK.

\subsubsubsection{Ferries}
Apart from the cargo ships, steamboats for tourists and passenger boats there are also a few ferries operating under the PID. 
The students of Matfyz often use the P7 ferry from Holešovice to Karlín which will take you very close to the MFF's building in Karlín.
Another important ferry is the P8 ferry that substitutes the collapsed Trojská lávka (a pedestrian bridge).
Ferries do not operate during winters and when the flood risk is increased.

\subsubsubsection{Night trams and buses}
Night trams operate approximately from 11:45 p.m. to 5:00 a.m. The routes are different than those of daytime trams. 
All of them meet at Lazarská stop (by the Spálená street, not far from Národní třída) where there is always enough time to change from one line to another; but there are of course more meet up points like this.
The intervals between connections are 30 minutes on weekdays, 20 minutes on weekends. Therefore it is necessary to pay attention and get off at the correct stop.
A lots of people who fall asleep find themselves right at the opposite side of Prague where the driver kicks them off the tram and lets them freeze.
The night tram network is complemented by the network of night buses which operate under a bit different rules (e.g. they do not meet and some of them run at hourly intervals).
Because there is almost no traffic at night some of the night bus lines tend to be more convenient than the daytime ones. 
For example travelling by the night line from Jižák to the student residence 17. listopadu is much faster than if you travelled during the daytime.
The most convenient bus stop to get from the city center at night is I. P. Pavlova where either the buses going in the direction of Kuchyňka and Volha stop.

The night tram numbers range from 91 to 99, the night buses have numbers ranging from 901 to 915 and the suburban night buses' numbers range from 951 to 960.

\subsubsubsection{Transit restrictions}
Traffic closures and transit restrictions are very common in Prague. The information about them appears at the bus and tram stops and in the metro stations and also buses and trams with uncommon numbers start running.
Trams substituting the common lines usually have numbers ranging from 31 to 49, buses are usually marked with the letter X followed by a number (which usually makes sense so if it is e.g. necessary to substitute the tram no. 3, the bus will be named X3). 
The substitutions for metro are usually named X(A,B,C).

\subsubsubsection{Bike}
Its main advantage is that it is faster then bus or tram on most of the routes. 
If you need to travel longer distances, you can bring your bike with you onto all kinds of public transport with the exception of buses.
You do not need to pay anything for bicycle transportation if you have a ticket for yourself.
In some of the metro stations you can use a lift but most of the times you will have to carry your bike with you on the escalator.
While travelling by metro the bike can be transported on each of the first and the last platforms of the individual carriages with the exception of the train’s first platform. It is best to use the last carriage.

\subsubsubsection{Kick Scooter}
The folding scooter is great for travelling shorter distances (e.g. from the dormitory to the metro) and you can transport it in all types of public transport for free even if the bus drivers can sometimes be reluctant to do so.
The kick scooter is, by the rules of public transport, considered as bike.

\subsubsubsection{Transport Sharing}
Using transport sharing (such as sharing bikes or kick scooters) in Prague is recently becoming more and more popular.
The most popular company providing this service is \textit{Rekola}. They lend bikes (including electric bikes) and all of them are pink. You can also borrow electrict bikes (yellow this time) from the company \textit{Freebike}. Electric scooters can be lent from the company \textit{Lime}.


The main advantage of using transport sharing is that you usually have a bike or a scooter available nearby and you do not have to deal with parking it. You just arrive to your destination and park next to a lamp.
Also this type of transport is usually considerably faster than using the public transport when travelling short distances. 
You just need to take into consideration the fact that there are some places where the bikes and scooters cannot be parked.


You can rent a bike or a scooter via smartphone apps created by the companies that offer this service. The prices vary a lot and usually depend on the time spent riding their bikes and scooters.


\subsubsection{Transport to the student residence 17. listopadu}
\subsubsubsection{How to get to there}
The most important access point is the metro station Nádraží Holešovice (``the Holešovice railway station"). From there you can either walk or go by bus. 
It takes fifteen minutes to walk to the dormitory but if you \textit{really} hurry, it is possible to make it in seven minutes.
If you want to go by bus, you can either take the line no. 112 and get off at Pelc Tyrolka (which is the closest bus stop to the dormitory) or take the line no. 201 and get off at Kuchyňka. 
Before making a decision which bus to choose, you should probably know that an average students needs 3 or 4 minutes more to walk to the dormitory from Kuchyňka than from Pelc Tyrolka.
You cannot travel by the bus no. 112 in the opposite direction (meaning from the dormitory) as it does not go this way.

You can use the tram no. 17 when travelling from somewhere around Výstaviště (e.g. from Strossmayerák or Právnická fakulta) and get off at Trojská (if you are travelling from the dormitory, you need to follow the road to the zoo until you run across the tramway).
It takes about nine minutes to get to this tram stop; another option includes using the bus no. 112 in the direction of the zoo at the stop Povltavská.

Every student of Matfyz sooner or later finds out that there are two tram stops named ``Nádraží Holešovice", usually at the moment he gets off at the wrong one and tries to figure out where the hell he is now. 
The tram no. 17 stops at the further one of them.

\subsubsection{Ideal Public Transport Algorithms}
Below you will find a few recommended connections and routes that might become useful for you. Of course if you have some spare time on your hands and are not easily intimidated, do not be afraid to try and find your own best route.

Following abbreviations are used while describing the recommended routes:

\begin{tabularx}{\textwidth}{ |l|X| }
\hline
M C (X \ra Y) & metro, line C, from~X to Y \\
\hline
T 8, 24 (X \ra Y)/2 & trams no. 8 a 24 from~X to Y and it is only two stops \\
\hline
B 201 (X \ra Y)/2 & bus no. 201 (the rest is the same as for trams)\\
\hline
V~221 (X \ra Y) & train on the route no. 221 from~X to Y \\
\hline
\end{tabularx}

\subsubsubsection{Routes from the student residence 17. listopadu to the faculty buildings}
Your route from the dormitory will most of the times take you via Nádraží Holešovice and there are two ways to get there so you may use the following substitution: B 201
(Kuchyňka \ra Nádrhol)/2 = B 112 (Pelc Tyrolka \ra Povltavská)/1, T 17 (Trojská
\ra Nádrhol (sever))/1 = B 112 (Pelc Tyrolka \ra Povltavská)/1, B 112 (Trojská
\ra Nádrhol (jih))/1. 
Another option includes walking across the bridge and taking the first right in the direction of Nádrhol. Besides the fact that it takes a really short time to get there and it is very reliable, you reach the metro from the side that will be the closest to the exit when you will later get off at Pavlák/Florenc.
If you do not need to get to the metro but you want to travel in the direction of Strossmayerák, it is worth it to walk to Trojská (it takes less than 10 minutes).

\textbf{Troja (physics - lecture rooms T)}\\
Walk in the direction of the bridge and go underneath it, then enter the low rise blue-gray building joined with the high rise building of the same colour.

\textbf{Troja (English - lecture rooms V)}\\
Walk in the direction of the bridge and before you get underneath it,  bypass the parking lot-junkyard from the left, then go under the bridge and straight to the low rise building which is not joined with the high rise building of the same colour.

\textbf{Karlín}\\
B 201 (Kuchyňka \ra Nádrhol)/2, M C (Nádrhol \ra Florenc), T 3, 8, 24 (Florenc
\ra Křižíkova)/2, cross the street and enter the building right in front of you or you can also walk from Florenc. This choice is the best on average but is a bit complicated.

The alternative includes only one change at Bulovka, tj. B 201
(Kuchyňka \ra Bulovka)/3, T 3, 24 (Bulovka \ra Křižíkova)/8. The building is on the right side, just a few tens of meters away from the tram stop.

Riding bike alongside the river is also fast: follow the cycle route A2 from the track in the direction of Palmovka through Thomayerovy sady and Rohanský ostrov to Karlín. 
There is a bicycle stand in the yard of the MFF's building. You can get to the yard through the second door, ring the reception.

\textbf{Karlov} \\
B 201 (Kuchyňka \ra Nádrhol)/2, M C (Nádrhol \ra Pavlák), get out in the direction of the street Na Bojišti, cross the street at the traffice lights, walk on the right side in the direction of Vyšehrad for a while, turn right (street Na Bojišti), then turn left at the end and continue straight ahead. The last two building on the right side are the buildings Ke Karlovu 3 and 5 (in the opposite order).
It is also possible to go one stop by bus no. 148 from Pavlák to the stop Dětská nemocnice Karlov and closer to the building this way.

\textbf{Malá Strana}\\
There are a few ways to get there by bus, metro or tram.
Probably the fastest one is to use the bus no. 112 or walk to Trojská and then T 17 (Trojská \ra Čechův most)/6, T 15 (Čechův most \ra
Malostranské náměstí)/2.

Depending on the current time of the day and season of the year it is almost equally fast (or a bit slower) to use the metro, which means B 201 (Kuchyňka \ra Nádrhol)/2 (or walk there), M C (Nádrhol \ra Muzeum), M A (Muzeum \ra Malostranská), T 12, 15, 20,
22, 23 (Malostranská \ra Malostranské náměstí)/1.

If you do not like to change and are not in a hurry, you can get to Malá strana by the tram no. 12 from Holešovice. Then we recommend to take a book with you and maybe even a snack.

Because the trafic near Malostranská is often heavy, it might be sometimes better to go from (or more often to Malostranská) on foot. You can pass through Valdštejnská zahrada (``Valdštejn garden") on the way.

You can also ride a bike:
People who do not mind riding a bike in a heavy traffic can go from the dormitory straight ahead, follow the road Magistrála to Nábřeží Kapitána Jaroše and then continue along the river. This way you can get to Malostranská in just 15 minutes.
Or you can follow the bicycle route in the direction of the zoo, use the ferry substituting Trojská lávka (which operates at the frequency of 10 min), go to Stromovka and up to Letenský park. 
Finally go through Chotkovy sady and the route on the right before Malostranská will lead you to the square right in front of the school.
You can leave your bike on the right side of the corridor (next to the vending machines) or in the yard.

\textbf{Hostivař}\\
B 201 (Kuchyňka \ra Nádrhol)/2, M C (Nádrhol \ra Hlavák), V 221 (Hlavák \ra
Hostivař) -- the train operates under the name S9, then walk on the street alongside the railway in the direction of Benešov until SCUK emerges from behind the blocks of flats. Or walk around the pond and then to SCUK through the housing estate.
The alternative is to walk across the street, then down to the terminal tram station and B
125 (Nádr. Hostivař \ra Gercenova)/1.

Or B 201 (Kuchyňka \ra Nádrhol)/2, M C (Nádrhol \ra Muzeum), M A (Muzeum \ra
Skalka), B 125 (Skalka \ra Gercenova)/7, take the diagonal street (Gercenova) and bypass the shopping center from the left.

Or: B 201 (Kuchyňka \ra Střížkov)/7, B 183 (Střížkov \ra Gercenova)/18. This option is reliable only outside rush hour.

\subsubsubsection{Routes between the faculty buildings}
\textbf{Karlov --- Karlín}\\
Walk to Pavlák, M C (Pavlák \ra Florenc), T 3, 8, 24 (Florenc \ra Křižíkova)/2, cross the street and enter the building.

\textbf{Troja --- wherever}\\
the same as if travelling from the dormitory

\textbf{Karlov --- Malá Strana}\\
Walk down the street Ke Karlovu, after you reach the tramway go left and after the next crossroad there is the tram stop Štěpánská on the right side, T 22, 23 (Štěpánská \ra
Malostranské náměstí)/6. You can shorten your route by going through Kateřinská zahrada (``Kateřinská garden"). 
If you are afraid to go to Štěpánská, you can walk to Pavlák and T 22, 23 (Pavlák \ra
Malostranské náměstí)/7. 

\textbf{Karlín --- Malá Strana}\\
M B (Křižíkova \ra Národní třída), T 22, 23 (Národní třída \ra Malostranské
náměstí)/4.

Or: T 8 (Křižíkova \ra Náměstí Republiky)/4, T 15 (Náměstí Republiky \ra
Malostranské náměstí)/4.

Or: M B (Křižíkova \ra Můstek), M A (Můstek \ra Malostranská), T 12, 15, 20,
22, 23 (Malostranská \ra Malostranské náměstí)/1. This option is often faster than those above if you really are in a hurry (you just run down the stairs, run through the station in the direction of the change to metro A, and run upstairs). 
This is the best way to take if you travel between these two buildings during the ten minute break (but you will still arrive to the next lesson about 10 min late).